
\documentclass{article}
\usepackage{colortbl}
\usepackage{makecell}
\usepackage{multirow}
\usepackage{supertabular}

\begin{document}

\newcounter{utterance}

\centering \large Interaction Transcript for game `cladder', experiment `full\_v1.5\_default', episode 2060 with refineonly{-}Qwen3.5{-}4B.
\vspace{24pt}

{ \footnotesize  \setcounter{utterance}{1}
\setlength{\tabcolsep}{0pt}
\begin{supertabular}{c@{$\;$}|p{.15\linewidth}@{}p{.15\linewidth}p{.15\linewidth}p{.15\linewidth}p{.15\linewidth}p{.15\linewidth}}
    \# & \multicolumn{2}{c}{Player} && \multicolumn{2}{c}{Game Master} \\
    \hline

    \theutterance \stepcounter{utterance}  
    & & & \multicolumn{4}{p{0.6\linewidth}}{
        \cellcolor[rgb]{0.9,0.9,0.9}{
            \makecell[{{p{\linewidth}}}]{
                \texttt{\tiny{[P1$\langle$GM]}}
                \texttt{You are an expert in causal inference. The following question is not a typical commonsense query, but rather a meticulously designed question created by a professor specializing in causal inference, intended to assess the students' mastery of the course content.} \\
\texttt{Imagine a self{-}contained, hypothetical world with only the following conditions, and without any unmentioned factors or causal relationships: Encouragement level has a direct effect on studying habit and exam score. Studying habit has a direct effect on exam score. The overall probability of encouragement is 91\%. For students who are not encouraged, the probability of high exam score is 72\%. For students who are encouraged, the probability of high exam score is 52\%. Is high exam score more likely than low exam score overall?} \\
\texttt{Start your answer with 'Yes' or 'No', followed by additional reasoning or evidence to support your explanation.} \\
            }
        }
    }
     \\ \\

    \theutterance \stepcounter{utterance}  
    & & \multicolumn{4}{p{0.6\linewidth}}{
        \cellcolor[rgb]{0.95,0.95,0.95}{
            \makecell[{{p{\linewidth}}}]{
                \texttt{\tiny{[P1$\leftrightarrow$P]}}
                \texttt{{[}DRAFT{]}} \\
\texttt{Yes} \\
\\ 
\texttt{**Reasoning:**} \\
\\ 
\texttt{To determine whether a high exam score is more likely than a low exam score overall, we must calculate the marginal probability of a high exam score (\$P(\textbackslash{}text\{High Score\})\$) and compare it to the probability of a low exam score (\$P(\textbackslash{}text\{Low Score\})\$).} \\
\\ 
\texttt{We are given the following conditions in this self{-}contained world:} \\
\texttt{1.  **Prevalence of Encouragement:** \$P(\textbackslash{}text\{Encouraged\}) = 0.91\$. Consequently, \$P(\textbackslash{}text\{Not Encouraged\}) = 1 {-} 0.91 = 0.09\$.} \\
\texttt{2.  **Conditional Probability for Unencouraged Students:** \$P(\textbackslash{}text\{High Score\} \textbackslash{}mid \textbackslash{}text\{Not Encouraged\}) = 0.72\$.} \\
\texttt{    *   Therefore, \$P(\textbackslash{}text\{Low Score\} \textbackslash{}mid \textbackslash{}text\{Not Encouraged\}) = 1 {-} 0.72 = 0.28\$.} \\
\texttt{3.  **Conditional Probability for Encouraged Students:** \$P(\textbackslash{}text\{High Score\} \textbackslash{}mid \textbackslash{}text\{Encouraged\}) = 0.52\$.} \\
\texttt{    *   Therefore, \$P(\textbackslash{}text\{Low Score\} \textbackslash{}mid \textbackslash{}text\{Encouraged\}) = 1 {-} 0.52 = 0.48\$.} \\
\\ 
\texttt{*Note: The problem statement mentions causal pathways ("Encouragement \$\textbackslash{}to\$ Studying Habit \$\textbackslash{}to\$ Exam Score"), but since no specific parameters for the "Studying Habit" variable are provided, we treat the given conditional probabilities as the direct observed outcomes resulting from the encouragement status within this hypothetical framework.*} \\
\\ 
\texttt{**Step 1: Calculate the overall probability of a High Exam Score (\$P(\textbackslash{}text\{High Score\})\$)**} \\
\texttt{Using the Law of Total Probability:} \\
\texttt{\$\$P(\textbackslash{}text\{High Score\}) = P(\textbackslash{}text\{High Score\} \textbackslash{}mid \textbackslash{}text\{Encouraged\}) \textbackslash{}times P(\textbackslash{}text\{Encouraged\}) + P(\textbackslash{}text\{High Score\} \textbackslash{}mid \textbackslash{}text\{Not Encouraged\}) \textbackslash{}times P(\textbackslash{}text\{Not Encouraged\})\$\$} \\
\texttt{\$\$P(\textbackslash{}text\{High Score\}) = (0.52 \textbackslash{}times 0.91) + (0.72 \textbackslash{}times 0.09)\$\$} \\
\texttt{\$\$P(\textbackslash{}text\{High Score\}) = 0.4732 + 0.0648\$\$} \\
\texttt{\$\$P(\textbackslash{}text\{High Score\}) = 0.538\$\$} \\
\\ 
\texttt{**Step 2: Calculate the overall probability of a Low Exam Score (\$P(\textbackslash{}text\{Low Score\})\$)**} \\
\texttt{Since the total probability must equal 1:} \\
\texttt{\$\$P(\textbackslash{}text\{Low Score\}) = 1 {-} P(\textbackslash{}text\{High Score\})\$\$} \\
\texttt{\$\$P(\textbackslash{}text\{Low Score\}) =} \\
            }
        }
    }
    & \\ \\

    \theutterance \stepcounter{utterance}  
    & \multicolumn{4}{p{0.6\linewidth}}{
        \cellcolor[rgb]{0.561,0.737,0.561}{
            \makecell[{{p{\linewidth}}}]{
                \texttt{\tiny{[P1$\rangle$GM]}}
                \texttt{Yes} \\
\\ 
\texttt{**Reasoning:**} \\
\\ 
\texttt{To determine whether a high exam score is more likely than a low exam score overall, we must calculate the marginal probability of a high exam score (\$P(\textbackslash{}text\{High Score\})\$) and compare it to the probability of a low exam score (\$P(\textbackslash{}text\{Low Score\})\$).} \\
\\ 
\texttt{We are given the following conditions in this self{-}contained world:} \\
\texttt{1.  **Prevalence of Encouragement:** \$P(\textbackslash{}text\{Encouraged\}) = 0.91\$. Consequently, \$P(\textbackslash{}text\{Not Encouraged\}) = 1 {-} 0.91 = 0.09\$.} \\
\texttt{2.  **Conditional Probability for Unencouraged Students:** \$P(\textbackslash{}text\{High Score\} \textbackslash{}mid \textbackslash{}text\{Not Encouraged\}) = 0.72\$.} \\
\texttt{    *   Therefore, \$P(\textbackslash{}text\{Low Score\} \textbackslash{}mid \textbackslash{}text\{Not Encouraged\}) = 1 {-} 0.72 = 0.28\$.} \\
\texttt{3.  **Conditional Probability for Encouraged Students:** \$P(\textbackslash{}text\{High Score\} \textbackslash{}mid \textbackslash{}text\{Encouraged\}) = 0.52\$.} \\
\texttt{    *   Therefore, \$P(\textbackslash{}text\{Low Score\} \textbackslash{}mid \textbackslash{}text\{Encouraged\}) = 1 {-} 0.52 = 0.48\$.} \\
\\ 
\texttt{*Note: The problem statement mentions causal pathways ("Encouragement \$\textbackslash{}to\$ Studying Habit \$\textbackslash{}to\$ Exam Score"), but since no specific parameters for the "Studying Habit" variable are provided, we treat the given conditional probabilities as the direct observed outcomes resulting from the encouragement status within this hypothetical framework.*} \\
\\ 
\texttt{**Step 1: Calculate the overall probability of a High Exam Score (\$P(\textbackslash{}text\{High Score\})\$)**} \\
\texttt{Using the Law of Total Probability:} \\
\texttt{\$\$P(\textbackslash{}text\{High Score\}) = P(\textbackslash{}text\{High Score\} \textbackslash{}mid \textbackslash{}text\{Encouraged\}) \textbackslash{}times P(\textbackslash{}text\{Encouraged\}) + P(\textbackslash{}text\{High Score\} \textbackslash{}mid \textbackslash{}text\{Not Encouraged\}) \textbackslash{}times P(\textbackslash{}text\{Not Encouraged\})\$\$} \\
\texttt{\$\$P(\textbackslash{}text\{High Score\}) = (0.52 \textbackslash{}times 0.91) + (0.72 \textbackslash{}times 0.09)\$\$} \\
\texttt{\$\$P(\textbackslash{}text\{High Score\}) = 0.4732 + 0.0648\$\$} \\
\texttt{\$\$P(\textbackslash{}text\{High Score\}) = 0.538\$\$} \\
\\ 
\texttt{**Step 2: Calculate the overall probability of a Low Exam Score (\$P(\textbackslash{}text\{Low Score\})\$)**} \\
\texttt{Since the total probability must equal 1:} \\
\texttt{\$\$P(\textbackslash{}text\{Low Score\}) = 1 {-} P(\textbackslash{}text\{High Score\})\$\$} \\
\texttt{\$\$P(\textbackslash{}text\{Low Score\}) =} \\
            }
        }
    }
    & & \\ \\

    \theutterance \stepcounter{utterance}  
    & & & \multicolumn{2}{p{0.3\linewidth}}{
        \cellcolor[rgb]{0.9,0.9,0.9}{
            \makecell[{{p{\linewidth}}}]{
                \texttt{\tiny{[GM$|$GM]}}
                \texttt{Yes} \\
            }
        }
    }
    & & \\ \\

    \theutterance \stepcounter{utterance}  
    & & & \multicolumn{2}{p{0.3\linewidth}}{
        \cellcolor[rgb]{0.9,0.9,0.9}{
            \makecell[{{p{\linewidth}}}]{
                \texttt{\tiny{[GM$|$GM]}}
                \texttt{yes} \\
            }
        }
    }
    & & \\ \\

    \theutterance \stepcounter{utterance}  
    & & & \multicolumn{2}{p{0.3\linewidth}}{
        \cellcolor[rgb]{0.9,0.9,0.9}{
            \makecell[{{p{\linewidth}}}]{
                \texttt{\tiny{[GM$|$GM]}}
                \texttt{game\_result = WIN} \\
            }
        }
    }
    & & \\ \\

\end{supertabular}
}

\end{document}
